% Luc Destombe --- licence CC BY-NC-SA 4.0
% https://creativecommons.org/licenses/by-nc-sa/4.0/deed.fr
% Fichier autonome : compiler avec pdflatex, deux fois (signets, nombre de pages).
\documentclass[11pt,a4paper]{article}
\makeatletter
\newif\ifeleve
\elevefalse

\newif\iflycee@palatino
\lycee@palatinotrue

\RequirePackage[utf8]{inputenc}
\RequirePackage[T1]{fontenc}
\RequirePackage[french]{babel}
\RequirePackage{etoolbox}

\newcommand{\ShorthandsOff}{\@ifundefined{shorthandoff}{}{\shorthandoff{;:!?}}}
\newcommand{\ShorthandsOn}{\@ifundefined{shorthandon}{}{\shorthandon{;:!?}}}
\ShorthandsOff
\AtBeginDocument{\ShorthandsOn}
\AtBeginEnvironment{tikzpicture}{\ShorthandsOff}

\RequirePackage{geometry}
\RequirePackage{fancyhdr}
\RequirePackage{lastpage}
\RequirePackage{multicol}
\RequirePackage{array}
\RequirePackage{enumitem}
\RequirePackage{xcolor}
\RequirePackage{needspace}

\RequirePackage{amsmath,amssymb}
\RequirePackage{mathpazo}
\DeclareMathSizes{10.95}{10.95}{8}{5.5}
\begingroup\catcode`\_=\active \catcode`\^=\active
\gdef_#1{\sb{\scriptscriptstyle #1}}
\gdef^#1{\sp{\scriptscriptstyle\lycee@moinsepais #1}}
\endgroup
\newcommand\lycee@moins{\mathbin{%
  \setbox\z@\hbox{$\m@th\scriptscriptstyle\lycee@moinsorig$}%
  \dimen@=.55\wd\z@ \advance\dimen@-.5pt
  \hbox to.75\wd\z@{\hss$\m@th\scriptscriptstyle\vcenter{\hbox{%
    \tikz\draw[line width=.5pt,line cap=round](0,0)--(\the\dimen@,0);}}$\hss}}}
\begingroup\lccode`\~=`\- \lowercase{\endgroup\let~}\lycee@moins
\newcommand\lycee@moinsepais{\mathcode`\-="8000 }
\AtBeginDocument{\mathchardef\lycee@moinsorig=\mathcode`\-\relax}
\AtBeginDocument{\catcode`\_=12 \mathcode`\_="8000
  \catcode`\^=12 \mathcode`\^="8000 }
\newcommand\lycee@exposants{%
  \fontdimen13\textfont2=.48\fontdimen6\textfont2
  \fontdimen14\textfont2=.48\fontdimen6\textfont2
  \fontdimen15\textfont2=.48\fontdimen6\textfont2
  \fontdimen13\scriptfont2=.48\fontdimen6\scriptfont2
  \fontdimen14\scriptfont2=.48\fontdimen6\scriptfont2
  \fontdimen15\scriptfont2=.48\fontdimen6\scriptfont2}
\every@math@size\expandafter{\the\every@math@size\lycee@exposants}
\newlength{\retraitcalculs}
\setlength{\retraitcalculs}{2em}

\RequirePackage{tikz}
\usetikzlibrary{arrows.meta,calc,babel,shapes.misc}
\RequirePackage[most]{tcolorbox}
\tcbuselibrary{skins,breakable}

\definecolor{coulDef}{HTML}{1F4E79}
\definecolor{coulProp}{HTML}{7030A0}
\definecolor{coulRem}{HTML}{C55A11}
\definecolor{coulEx}{HTML}{2E7D32}
\definecolor{coulExo}{HTML}{B03A2E}
\definecolor{coulPart}{HTML}{1F4E79}
\definecolor{coulMeth}{HTML}{1B6E4F}
\definecolor{coulCourbe}{HTML}{0B5ED7}
\definecolor{coulCourbeB}{HTML}{C00000}
\definecolor{coulCourbeC}{HTML}{2E7D32}
\definecolor{coulGrille}{HTML}{8C8C8C}
\definecolor{coulRep}{HTML}{1B6E4F}
\definecolor{coulBar}{HTML}{2E7D32}

\newcommand{\R}{\mathbb{R}}
\newcommand{\Rs}{\mathbb{R}^{*}}
\renewcommand{\le}{\leqslant}
\renewcommand{\ge}{\geqslant}
\newcommand{\vect}[1]{\overrightarrow{#1}}
\newcommand{\norme}[1]{\left\lVert #1 \right\rVert}
\newcommand{\intervalle}[2]{\left[#1\,;#2\right]}
\RequirePackage{cancel}
\renewcommand{\CancelColor}{\color{coulExo}}
\newcommand{\fois}[1]{\times{\color{coulExo}#1}}
\newcommand{\barre}[1]{\cancel{#1}}

\tikzset{grille/.style={coulGrille,line width=0.4pt}}
\newcommand{\quadrillage}[4]{\draw[grille,step=1] (#1,#2) grid (#3,#4);}
\tikzset{
  croix/.style={cross out,draw,line width=0.6pt,inner sep=0pt,minimum size=4pt},
  croixdroite/.style={croix,rotate=45,line width=0.8pt},
}
\newcommand{\pt}[3]{\path (#1) node[croix]{} node[#2,font=\small,inner sep=2.5pt]{$#3$};}

\newcommand{\lycee@repere}[4]{%
  \draw[grille,step=1] (#1,#2) grid (#3,#4);
  \draw[-{Stealth[length=2mm]},thick] (#1,0)--({#3+0.4},0) node[below left=-1pt]{$x$};
  \draw[-{Stealth[length=2mm]},thick] (0,#2)--(0,{#4+0.4}) node[below left=-1pt]{$y$};
  \node[below left,font=\scriptsize,inner sep=1.5pt] at (0,0) {$0$};}
\newcommand{\lycee@gradx}[1]{\foreach \k/\l in {#1}{%
  \draw (\k,0.07)--(\k,-0.07) node[below,font=\scriptsize,inner sep=1.5pt]{$\l$};}}
\newcommand{\lycee@grady}[1]{\foreach \k/\l in {#1}{%
  \draw (0.07,\k)--(-0.07,\k) node[left,font=\scriptsize,inner sep=1.5pt]{$\l$};}}
\newcommand{\lycee@intff}[2]{\left[#1\,;#2\right]}
\newcommand{\lycee@intfo}[2]{\left[#1\,;#2\right[}
\newcommand{\lycee@intof}[2]{\left]#1\,;#2\right]}
\newcommand{\lycee@intoo}[2]{\left]#1\,;#2\right[}
\newcommand{\lycee@ens}[1]{\left\{#1\right\}}
\AtBeginDocument{%
  \providecommand{\repere}{\lycee@repere}%
  \providecommand{\gradx}{\lycee@gradx}%
  \providecommand{\grady}{\lycee@grady}%
  \providecommand{\Cf}{\mathcal{C}\sb{\scriptscriptstyle f}}%
  \providecommand{\Cg}{\mathcal{C}\sb{\scriptscriptstyle g}}%
  \providecommand{\intff}{\lycee@intff}%
  \providecommand{\intfo}{\lycee@intfo}%
  \providecommand{\intof}{\lycee@intof}%
  \providecommand{\intoo}{\lycee@intoo}%
  \providecommand{\ens}{\lycee@ens}}

\newlist{questions}{enumerate}{3}
\setlist[questions,1]{label=\arabic*\textdegree),leftmargin=*,itemsep=2pt,topsep=3pt}
\setlist[questions,2]{label=\alph*),leftmargin=*,itemsep=1pt,topsep=2pt}
\setlist[questions,3]{label=\roman*),leftmargin=*,itemsep=1pt,topsep=1pt}
\setlist[itemize]{label=\textbullet,leftmargin=*,itemsep=2pt,topsep=3pt}

\newcommand{\besoinplace}[1]{\par\penalty\@M\begingroup
  \dimen@=#1\relax\dimen@ii=\pagegoal\advance\dimen@ii-\pagetotal
  \ifdim\dimen@>\dimen@ii\ifdim\dimen@ii>\z@\vfil\fi\break\fi\endgroup}

\newcommand{\lycee@pied}{\thepage/\pageref{LastPage}}
\newcommand{\entete}[2]{%
  \pagestyle{fancy}\fancyhf{}%
  \fancyhead[L]{\footnotesize\color{coulPart}\textbf{#1}}%
  \fancyhead[R]{\footnotesize\color{coulPart}\textbf{#2}}%
  \fancyfoot[C]{\footnotesize\lycee@pied}%
  \renewcommand{\headrulewidth}{0.4pt}%
  \renewcommand{\headrule}{\hbox to\headwidth{%
    \color{coulPart!40}\leaders\hrule height \headrulewidth\hfill}}}

\RequirePackage{ccicons}
\newcommand{\lycee@licence}{{\scriptsize\color{gray}%
  \copyright\ Luc Destombe --- \ccbyncsa\ CC BY-NC-SA 4.0}}
\newif\iflycee@licence \lycee@licencetrue
\newcommand{\sanslicence}{\lycee@licencefalse}
\AtBeginDocument{\iflycee@licence\AddToHookNext{shipout/foreground}{%
  \put(0,0){\rlap{\kern\dimexpr1in+\hoffset+\oddsidemargin+\textwidth\relax
    \llap{\raisebox{-\dimexpr1in+\voffset+\topmargin+\headheight+\headsep
      +\textheight+\footskip\relax}{\lycee@licence}}}}}\fi}

\newcommand{\formulesserrees}{%
  \setlength{\abovedisplayskip}{3pt}\setlength{\belowdisplayskip}{3pt}%
  \setlength{\abovedisplayshortskip}{2pt}\setlength{\belowdisplayshortskip}{3pt}}

\setlength{\parindent}{0pt}

\RequirePackage[implicit=false,hidelinks,pdfencoding=auto,psdextra,bookmarksopen,
  bookmarksopenlevel=1,pdfcreator={},pdfproducer={}]{hyperref}
\RequirePackage{bookmark}
\hypersetup{pdfauthor={Luc Destombe},pdfkeywords={CC BY-NC-SA 4.0}}
\newcounter{lycee@signet}
\newcommand{\signet}[3][]{\stepcounter{lycee@signet}%
  \edef\lycee@dest{\if\relax\detokenize{#1}\relax
    signet.\arabic{lycee@signet}\else#1\fi}%
  \hypertarget{\lycee@dest}{}\bookmark[level=#2,dest=\lycee@dest]{#3}}
\pdfstringdefDisableCommands{%
  \def\R{R}\def\vect#1{#1}\def\Cf{Cf}\def\Cg{Cg}\def\fois#1{×#1}%
  \def\barre#1{#1}\def\intervalle#1#2{[#1;#2]}\def\intff#1#2{[#1;#2]}%
  \def\textsuperscript#1{#1}\def\dfrac#1#2{#1/#2}\def\frac#1#2{#1/#2}%
  \def\vec#1{#1⃗}}%
\geometry{top=2cm,bottom=1cm,left=1cm,right=1cm,headheight=14pt,footskip=0.6cm}

\RequirePackage{pgfplots}
\pgfplotsset{compat=1.18}
\RequirePackage{tkz-tab}

\newcounter{partiecours}
\renewcommand{\thepartiecours}{\Roman{partiecours}}
\newcommand{\partiecours}[1]{%
  \refstepcounter{partiecours}%
  \Needspace*{8\baselineskip}%
  \bigskip
  \noindent\begin{tcolorbox}[enhanced,colback=coulPart,colframe=coulPart,
      boxrule=0pt,arc=2pt,left=8pt,right=8pt,top=4pt,bottom=4pt,
      before skip=6pt,after skip=10pt]
    \color{white}\bfseries\large \thepartiecours{} --- #1%
    \signet[partie-\arabic{partiecours}]{1}{\thepartiecours{} --- #1}
  \end{tcolorbox}%
  \addcontentsline{toc}{section}{\thepartiecours{} --- #1}%
}

\newtcolorbox{boitecours}[3][]{%
  enhanced,breakable,
  colback=white,colframe=#2,coltitle=white,
  fonttitle=\bfseries,
  attach boxed title to top left={xshift=8pt,yshift=-\tcboxedtitleheight/2},
  boxed title style={colback=#2,arc=2pt,boxrule=0pt},
  boxrule=0.9pt,arc=2pt,
  left=8pt,right=8pt,top=8pt,bottom=6pt,
  title={#3},#1}

\newcounter{methode}
\newenvironment{definitionbox}[1][Définition]%
  {\begin{boitecours}[phantom={\signet{2}{#1}}]{coulDef}{#1}}{\end{boitecours}}
\newenvironment{proprietebox}[1][Propriété]%
  {\begin{boitecours}[phantom={\signet{2}{#1}}]{coulProp}{#1}}{\end{boitecours}}
\newenvironment{methodebox}[1][Méthode]{\stepcounter{methode}%
  \begin{boitecours}[phantom={\signet[methode-\arabic{methode}]{2}{#1}}]{coulMeth}{#1}}%
  {\end{boitecours}}

\newcommand{\etape}[1]{\par\smallskip\textbf{\color{coulMeth}#1}\par\nobreak\smallskip}
\newcommand{\enonce}[1]{\emph{#1}\par\smallskip}
\newcommand{\reponse}[1]{\par\smallskip
  {\footnotesize\itshape\color{black!55}Réponses : #1}}
\newcommand{\demo}[1]{\textbf{\color{coulMeth}Démonstration.} #1}
\newcommand{\afaire}[1]{%
  \par\smallskip
  \noindent{\small\color{coulExo}$\blacktriangleright$\ \textbf{À faire :}
    fiche d'exercices du chapitre, #1.}\par\smallskip}

\newenvironment{calculs}{\setlength{\jot}{5pt}\@fleqntrue\@mathmargin\retraitcalculs
  \setlength{\abovedisplayskip}{4pt}\setlength{\belowdisplayskip}{4pt}%
  \setlength{\abovedisplayshortskip}{4pt}\setlength{\belowdisplayshortskip}{4pt}%
  \csname alignat*\endcsname{2}}%
  {\csname endalignat*\endcsname}
\newcommand{\rem}[1]{\text{\small\itshape\color{black!60}#1}}
\newcommand{\explique}[2]{\par\smallskip\noindent
  \begin{minipage}[t]{0.57\linewidth}\raggedright #1\end{minipage}\hfill
  \begin{minipage}[t]{0.40\linewidth}\raggedright\leavevmode
    {\small\itshape\color{black!60}#2}\end{minipage}\par}
\newcommand{\cas}[1]{\par\smallskip\noindent\textbf{\color{coulExo}#1}\nobreak}
\newcommand{\calc}[1]{\par\smallskip\textbf{\color{coulExo}#1}\par\nobreak}

\newtcolorbox{remarquebox}[1][Remarque]{%
  enhanced,breakable,colback=white,colframe=coulRem,
  boxrule=0pt,leftrule=2.5pt,arc=0pt,outer arc=0pt,
  left=8pt,right=6pt,top=5pt,bottom=5pt,
  before upper={\textbf{\color{coulRem}#1}\par\smallskip}}

\newtcolorbox{exemplebox}[1][Exemple]{%
  enhanced,breakable,colback=white,colframe=coulEx,
  boxrule=0pt,leftrule=2.5pt,arc=0pt,outer arc=0pt,
  left=8pt,right=6pt,top=5pt,bottom=5pt,
  before upper={\textbf{\color{coulEx}#1}\par\smallskip}}

\pgfplotsset{
  repere/.style={
    axis lines=middle,
    axis line style={-{Stealth[length=2.4mm]},thick},
    every axis x label/.style={at={(ticklabel* cs:1.0)},anchor=west},
    every axis y label/.style={at={(ticklabel* cs:1.0)},anchor=south},
    label style={font=\footnotesize},
    tick label style={font=\scriptsize},
    grid=both,
    major grid style={grille},
    minor tick num=0,
    tick align=inside,
    clip mode=individual,
  },
}
\tikzset{
  courbe/.style={coulCourbe,line width=1.1pt,smooth},
  courbeB/.style={coulCourbeB,line width=1.1pt,smooth},
  courbeC/.style={coulCourbeC,line width=1.1pt,smooth},
}

\newlist{etapes}{enumerate}{1}
\setlist[etapes]{label=\textbf{\arabic*.},leftmargin=*,itemsep=1pt,topsep=2pt}

\newlist{expressions}{enumerate}{1}
\setlist[expressions]{label={\Alph*\ $=$},leftmargin=*,itemsep=3pt,topsep=3pt}

\newcounter{exo}
\renewcommand{\theexo}{\ifnum\value{exo}<10 0\fi\arabic{exo}}
\newtcolorbox{exercice}[1][]{%
  enhanced,breakable,
  colback=white,colframe=coulExo,
  boxrule=0.9pt,arc=2pt,
  left=8pt,right=8pt,top=8pt,bottom=6pt,
  coltitle=white,fonttitle=\bfseries,
  attach boxed title to top left={xshift=8pt,yshift=-\tcboxedtitleheight/2},
  boxed title style={colback=coulExo,arc=2pt,boxrule=0pt},
  phantom={\signet{2}{Exercice \arabic{exo}}},
  title={Exercice \theexo\ifx\relax#1\relax\else{} --- #1\fi}}
\newenvironment{exo}[1][\relax]{\refstepcounter{exo}\begin{exercice}[#1]}{\end{exercice}}

  \newenvironment{correction}{\par}{\par}
  \newcommand{\autableau}{}
  \newcommand{\etapeexemples}{\etape{Exemples corrigés}}
  \newcommand{\etapetableau}[1]{\etape{#1}}
  \newenvironment{exempletableau}[1][Exemple]%
    {\begin{exemplebox}[#1]}{\end{exemplebox}}

\RequirePackage{comment}
\excludecomment{correctionavj}

\newcommand{\coursEntete}{}
\newcommand{\coursNiveau}{}
\newcommand{\coursTitre}{}
\newcommand{\coursSurTitre}{}
\newcommand{\coursSousTitre}{}

\newcommand{\enteteCours}[1]{\renewcommand{\coursEntete}{#1}}
\newcommand{\chapitrecours}[2]{\enteteCours{Chapitre #1 --- #2}}
\newcommand{\niveaucours}[1]{\renewcommand{\coursNiveau}{#1}}
\newcommand{\titrecours}[3][]{%
  \renewcommand{\coursTitre}{#2}%
  \renewcommand{\coursSurTitre}{#1}%
  \renewcommand{\coursSousTitre}{#3}}

\pagestyle{fancy}
\fancyhf{}
\fancyhead[L]{\footnotesize\color{coulPart}\textbf{\coursEntete}}
\fancyhead[R]{\footnotesize\color{coulPart}\textbf{\coursNiveau}}
\fancyfoot[C]{\footnotesize\thepage}
\renewcommand{\headrulewidth}{0.4pt}
\renewcommand{\headrule}{\hbox to\headwidth{\color{coulPart!40}\leaders\hrule height \headrulewidth\hfill}}

\setlength{\parskip}{2pt}

\AtBeginDocument{%
  \begin{center}
    \begin{tcolorbox}[enhanced,width=0.72\linewidth,colback=white,colframe=coulPart,
        boxrule=1.6pt,arc=3pt,halign=center,top=10pt,bottom=10pt]
      \ifx\coursSurTitre\empty
        {\Huge\bfseries\color{coulPart} \coursTitre}\\[4pt]
      \else
        {\Huge\bfseries\color{coulPart} \coursTitre}\\[3pt]
        {\Large\bfseries\color{coulPart} \coursSurTitre}\\[5pt]
      \fi
      {\normalsize \coursSousTitre}%
      
    \end{tcolorbox}
  \end{center}%
}
\makeatother

\chapitrecours{5}{Variations, signe et inéquations}
\niveaucours{Seconde}
\titrecours{VARIATIONS, SIGNE ET INÉQUATIONS}{Cours --- Seconde}

\begin{document}

\providecommand{\Ch}{\mathcal{C}\sb{\scriptscriptstyle h}}
\providecommand{\Ck}{\mathcal{C}\sb{\scriptscriptstyle k}}

\begin{remarquebox}[Comment utiliser ce cours]
Chaque partie commence par ce qu'il faut \textbf{savoir} (définitions, propriétés), puis
vient ce qu'il faut \textbf{savoir faire} : une méthode, avec des
\textbf{exemples corrigés} faits en classe, puis des exercices
\og\textbf{À vous de jouer}\fg{} à chercher seul.

\end{remarquebox}

\partiecours{Fonction croissante, fonction décroissante}

Une courbe se lit comme le profil d'une étape de montagne du Tour de France : en la
parcourant \textbf{de gauche à droite}, on monte ou on descend.

\begin{definitionbox}[Définition 1 --- Fonction croissante, décroissante, constante]
Soit $f$ une fonction définie sur un intervalle $I$.
\begin{itemize}
  \item $f$ est \textbf{croissante} sur $I$ si, quand $x$ augmente, $f(x)$ augmente :
        pour tous nombres $a$ et $b$ de $I$, si $a<b$ alors $f(a)\le f(b)$. Sa courbe
        \textbf{monte}.
  \item $f$ est \textbf{décroissante} sur $I$ si, quand $x$ augmente, $f(x)$ diminue :
        si $a<b$ alors $f(a)\ge f(b)$. Sa courbe \textbf{descend}.
  \item $f$ est \textbf{constante} sur $I$ si $f(x)$ garde toujours la même valeur.
        Sa courbe est horizontale.
\end{itemize}

\smallskip
\begin{minipage}{0.48\linewidth}\centering
\begin{tikzpicture}[x=1cm,y=1cm]
  \draw[-{Stealth[length=2mm]},thick] (0,0)--(3.6,0) node[below left=-1pt]{$x$};
  \draw[-{Stealth[length=2mm]},thick] (0,0)--(0,2.6) node[below left=-1pt]{$y$};
  \draw[coulCourbe,line width=1.1pt,domain=0.4:3.2,samples=30]
    plot (\x,{0.2*\x*\x+0.3});
  \draw[dashed,black!60] (1,0) node[below,font=\small]{$a$}--(1,0.5)
    --(0,0.5) node[left,font=\small]{$f(a)$};
  \draw[dashed,black!60] (2.6,0) node[below,font=\small]{$b$}--(2.6,1.652)
    --(0,1.652) node[left,font=\small]{$f(b)$};
  \path[coulCourbe] (1,0.5) node[croix]{} (2.6,1.652) node[croix]{};
\end{tikzpicture}

{\small croissante : $a<b$ et $f(a)\le f(b)$}
\end{minipage}\hfill
\begin{minipage}{0.48\linewidth}\centering
\begin{tikzpicture}[x=1cm,y=1cm]
  \draw[-{Stealth[length=2mm]},thick] (0,0)--(3.6,0) node[below left=-1pt]{$x$};
  \draw[-{Stealth[length=2mm]},thick] (0,0)--(0,2.6) node[below left=-1pt]{$y$};
  \draw[coulCourbeB,line width=1.1pt,domain=0.4:3.2,samples=30]
    plot (\x,{2.4-0.2*\x*\x});
  \draw[dashed,black!60] (1,0) node[below,font=\small]{$a$}--(1,2.2)
    --(0,2.2) node[left,font=\small]{$f(a)$};
  \draw[dashed,black!60] (2.6,0) node[below,font=\small]{$b$}--(2.6,1.048)
    --(0,1.048) node[left,font=\small]{$f(b)$};
  \path[coulCourbeB] (1,2.2) node[croix]{} (2.6,1.048) node[croix]{};
\end{tikzpicture}

{\small décroissante : $a<b$ et $f(a)\ge f(b)$}
\end{minipage}

\smallskip
\textbf{Une fonction croissante conserve l'ordre ; une fonction décroissante
renverse l'ordre.}

\textit{Exemple :} $2<7$. Si $f$ est croissante sur $\intff{0}{10}$, alors
$f(2)\le f(7)$ ; si $g$ est décroissante sur $\intff{0}{10}$, alors $g(2)\ge g(7)$.

\smallskip
Donner les \textbf{variations} de $f$, c'est dire sur quels intervalles elle est
croissante et sur lesquels elle est décroissante.
\end{definitionbox}

\begin{methodebox}[Méthode 1 --- Lire les variations d'une fonction sur sa courbe]
\etapeexemples
\begin{minipage}[c]{0.58\linewidth}\raggedright
La fonction $f$ est représentée ci-contre par sa courbe $\Cf$.
\begin{questions}[label=\alph*)]
  \item Sur quel intervalle la fonction $f$ est-elle définie ?
  \item Sur quels intervalles $f$ est-elle croissante ? décroissante ?
\end{questions}
\end{minipage}\hfill
\begin{minipage}[c]{0.38\linewidth}\centering
\begin{tikzpicture}[x=0.42cm,y=0.42cm]
  \repere{-5}{-3}{6}{5}
  \gradx{-4,-3,-2,-1,1,2,3,4,5} \grady{-2,-1,1,2,3,4}
  \draw[coulCourbe,line width=1.1pt]
    (-4,1) .. controls (-3,2.8) and (-2,4) .. (-1,4)
           .. controls (0,4) and (0.33,2) .. (1,1)
           .. controls (1.67,0) and (2.33,-2) .. (3,-2)
           .. controls (3.67,-2) and (4.33,0) .. (5,1);
  \path[coulCourbe] (-4,1) node[croix]{} (5,1) node[croix]{};
  \node[coulCourbe,font=\footnotesize] at (-3.6,3.2) {$\Cf$};
\end{tikzpicture}
\end{minipage}
\begin{correction}
\cas{a)}
\explique{$f$ est définie sur $\intff{-4}{5}$.}{la courbe commence en $x=-4$ et
s'arrête en $x=5$}
\cas{b)}
\explique{$f$ est croissante sur $\intff{-4}{-1}$.}{de $x=-4$ à $x=-1$, la courbe monte}
\explique{$f$ est décroissante sur $\intff{-1}{3}$.}{de $x=-1$ à $x=3$, elle descend}
\explique{$f$ est croissante sur $\intff{3}{5}$.}{puis elle remonte jusqu'à $x=5$}
\explique{Les intervalles s'écrivent avec des valeurs de $x$.}{on ne recopie pas les
hauteurs $4$ ou $-2$}
\end{correction}

\etape{À vous de jouer}
\begin{minipage}[c]{0.58\linewidth}\raggedright
La fonction $g$ est représentée ci-contre par sa courbe $\Cg$.
\begin{questions}[label=\alph*)]
  \item Sur quel intervalle la fonction $g$ est-elle définie ?
  \item Donner les variations de $g$.
\end{questions}
\reponse{a) $\intff{-3}{4}$ ; b) $g$ décroissante sur $\intff{-3}{0}$, croissante sur
$\intff{0}{2}$, décroissante sur $\intff{2}{4}$.}
\end{minipage}\hfill
\begin{minipage}[c]{0.38\linewidth}\centering
\begin{tikzpicture}[x=0.42cm,y=0.42cm]
  \repere{-4}{-2}{5}{4}
  \gradx{-3,-2,-1,1,2,3,4} \grady{-1,1,2,3}
  \draw[coulCourbeB,line width=1.1pt]
    (-3,3) .. controls (-2,1.4) and (-1,-1) .. (0,-1)
           .. controls (0.67,-1) and (1.33,2) .. (2,2)
           .. controls (2.67,2) and (3.33,0.87) .. (4,0);
  \path[coulCourbeB] (-3,3) node[croix]{} (4,0) node[croix]{};
  \node[coulCourbeB,font=\footnotesize] at (-2.4,3.4) {$\Cg$};
\end{tikzpicture}
\end{minipage}
\end{methodebox}

\afaire{exercices 1 à 6}

\partiecours{Maximum, minimum}

\begin{definitionbox}[Définition 2 --- Maximum, minimum]
Soit $f$ une fonction définie sur un intervalle $I$, et $a$ un nombre de $I$.
\begin{itemize}
  \item $f$ admet un \textbf{maximum} $M$ en $a$ si $f(a)=M$ et si, pour tout $x$ de $I$,
        $f(x)\le M$ : c'est la plus grande image, au point le plus haut de la courbe.
  \item $f$ admet un \textbf{minimum} $m$ en $a$ si $f(a)=m$ et si, pour tout $x$ de $I$,
        $f(x)\ge m$ : c'est la plus petite image, au point le plus bas.
\end{itemize}
Un maximum ou un minimum s'appelle un \textbf{extremum}. On dit qu'il est
\textbf{atteint} en $a$.
\end{definitionbox}

\begin{methodebox}[Méthode 2 --- Lire les extremums d'une fonction sur sa courbe]
\etapeexemples
\begin{minipage}[c]{0.58\linewidth}\raggedright
La fonction $h$, définie sur $\intff{-4}{3}$, est représentée ci-contre par sa
courbe $\Ch$.
\begin{questions}[label=\alph*)]
  \item Donner le maximum et le minimum de $h$ sur $\intff{-4}{3}$, et où ils sont
        atteints.
  \item Même question sur l'intervalle $\intff{1}{3}$ seulement.
\end{questions}
\end{minipage}\hfill
\begin{minipage}[c]{0.38\linewidth}\centering
\begin{tikzpicture}[x=0.42cm,y=0.42cm]
  \repere{-5}{-4}{4}{4}
  \gradx{-4,-3,-2,-1,1,2,3} \grady{-3,-2,-1,1,2,3}
  \draw[coulCourbe,line width=1.1pt]
    (-4,1) .. controls (-3.33,-0.67) and (-2.67,-3) .. (-2,-3)
           .. controls (-1,-3) and (0,3) .. (1,3)
           .. controls (1.67,3) and (2.33,2) .. (3,1);
  \path[coulCourbe] (-4,1) node[croix]{} (3,1) node[croix]{};
  \node[coulCourbe,font=\footnotesize] at (-3.4,2) {$\Ch$};
\end{tikzpicture}
\end{minipage}
\begin{correction}
\cas{a)}
\explique{Le maximum de $h$ est $3$, atteint en $x=1$.}{point le plus haut :
$(1\,;3)$}
\explique{Le minimum de $h$ est $-3$, atteint en $x=-2$.}{point le plus bas :
$(-2\,;-3)$}
\cas{b)}
\explique{Sur $\intff{1}{3}$, le maximum est $3$, atteint en $x=1$.}{on ne regarde
que la partie de la courbe entre $x=1$ et $x=3$}
\explique{Le minimum est $1$, atteint en $x=3$.}{un extremum peut être au bord de
l'intervalle}
\end{correction}

\etape{À vous de jouer}
Donner le maximum et le minimum de la fonction $g$ de la méthode 1 sur $\intff{-3}{4}$,
et où ils sont atteints.
\reponse{maximum $3$, atteint en $x=-3$ (au bord, plus haut que le sommet $(2\,;2)$) ;
minimum $-1$, atteint en $x=0$.}
\end{methodebox}

\afaire{exercices 7 à 9}

\partiecours{Tableau de variations}

\begin{definitionbox}[Définition 3 --- Tableau de variations]
Le \textbf{tableau de variations} d'une fonction résume ses variations :
\begin{itemize}
  \item sur la première ligne, les valeurs de $x$ : les bornes de l'intervalle de
        définition et les valeurs où le sens de variation change ;
  \item sur la seconde ligne, une flèche qui monte (fonction croissante) ou qui
        descend (fonction décroissante), avec les images au bout des flèches.
\end{itemize}
\end{definitionbox}

\begin{methodebox}[Méthode 3 --- Dresser le tableau de variations à partir de la courbe]
\etapeexemples
\begin{minipage}[c]{0.58\linewidth}\raggedright
La fonction $k$, définie sur $\intff{-3}{4}$, est représentée ci-contre par sa
courbe $\Ck$. Dresser son tableau de variations.
\end{minipage}\hfill
\begin{minipage}[c]{0.38\linewidth}\centering
\begin{tikzpicture}[x=0.42cm,y=0.42cm]
  \repere{-4}{-3}{5}{5}
  \gradx{-3,-2,-1,1,2,3,4} \grady{-2,-1,1,2,3,4}
  \draw[coulCourbe,line width=1.1pt]
    (-3,4) .. controls (-1.67,1.33) and (-0.33,-2) .. (1,-2)
           .. controls (2,-2) and (3,0.67) .. (4,2);
  \path[coulCourbe] (-3,4) node[croix]{} (4,2) node[croix]{};
  \node[coulCourbe,font=\footnotesize] at (-2,4.2) {$\Ck$};
\end{tikzpicture}
\end{minipage}
\begin{correction}
\cas{Résolution}\par\nopagebreak
\explique{Première ligne : $-3$, $1$ et $4$.}{bornes $-3$ et $4$ ; changement de sens
en $1$}
\explique{Seconde ligne : une flèche qui descend, puis une qui monte.}{décroissante sur
$\intff{-3}{1}$, croissante sur $\intff{1}{4}$}
\explique{Au bout des flèches : $k(-3)=4$, $k(1)=-2$ et $k(4)=2$.}{images lues sur la
courbe}
\begin{center}
\begin{tikzpicture}
  \tkzTabInit[lgt=1.6,espcl=2.4,deltacl=0.6]{$x$/1,$k(x)$/2}{$-3$,$1$,$4$}
  \tkzTabVar{+/$4$,-/$-2$,+/$2$}
\end{tikzpicture}
\end{center}
\end{correction}

\etape{À vous de jouer}
Dresser le tableau de variations de la fonction $g$ de la méthode 1.
\reponse{$x$ : $-3$, $0$, $2$, $4$ ; flèches : descend de $3$ à $-1$, monte jusqu'à $2$,
descend jusqu'à $0$.}
\end{methodebox}

\begin{methodebox}[Méthode 4 --- Exploiter un tableau de variations]
\etapeexemples
\begin{minipage}[c]{0.5\linewidth}\raggedright
La fonction $s$ a le tableau de variations ci-contre.
\begin{questions}[label=\alph*)]
  \item Comparer $s(-4)$ et $s(-2)$.
  \item Comparer $s(0)$ et $s(3)$.
  \item Donner le maximum et le minimum de $s$.
\end{questions}
\end{minipage}\hfill
\begin{minipage}[c]{0.46\linewidth}\centering
\begin{tikzpicture}[scale=0.9]
  \tkzTabInit[lgt=1.3,espcl=1.6,deltacl=0.6]{$x$/1,$s(x)$/2}{$-5$,$-1$,$4$,$6$}
  \tkzTabVar{-/$2$,+/$7$,-/$-3$,+/$0$}
\end{tikzpicture}
\end{minipage}
\begin{correction}
\cas{a)}
\explique{$-4$ et $-2$ sont dans $\intff{-5}{-1}$, où $s$ est croissante.}{on cherche
l'intervalle qui contient les deux nombres}
\explique{$-4<-2$, donc $s(-4)\le s(-2)$.}{croissante : l'ordre est conservé}
\cas{b)}
\explique{$0$ et $3$ sont dans $\intff{-1}{4}$, où $s$ est décroissante.}{même
intervalle}
\explique{$0<3$, donc $s(0)\ge s(3)$.}{décroissante : l'ordre est inversé}
\cas{c)}
\explique{Le maximum de $s$ est $7$, atteint en $x=-1$.}{plus grand nombre de la
seconde ligne}
\explique{Le minimum de $s$ est $-3$, atteint en $x=4$.}{plus petit nombre de la
seconde ligne}
\end{correction}

\etape{À vous de jouer}
\begin{minipage}[c]{0.5\linewidth}\raggedright
La fonction $t$ a le tableau de variations ci-contre.
\begin{questions}[label=\alph*)]
  \item Comparer $t(-2)$ et $t(0)$, puis $t(2)$ et $t(4)$.
  \item Donner le maximum et le minimum de $t$.
\end{questions}
\reponse{a) $t(-2)\le t(0)$ ; $t(2)\ge t(4)$. b) maximum $5$ en $x=1$ ; minimum $-1$
en $x=-3$.}
\end{minipage}\hfill
\begin{minipage}[c]{0.46\linewidth}\centering
\begin{tikzpicture}[scale=0.9]
  \tkzTabInit[lgt=1.3,espcl=2.2,deltacl=0.6]{$x$/1,$t(x)$/2}{$-3$,$1$,$5$}
  \tkzTabVar{-/$-1$,+/$5$,-/$2$}
\end{tikzpicture}
\end{minipage}
\end{methodebox}

\afaire{exercices 10 à 14}

\partiecours{Fonctions affines et linéaires}

\begin{definitionbox}[Définition 4 --- Fonction affine, fonction linéaire]
Une fonction $f$ est \textbf{affine} s'il existe deux nombres $a$ et $b$ tels que,
pour tout $x$ de $\R$ :
\[
  f(x)=ax+b
\]
\begin{itemize}
  \item Si $b=0$, $f(x)=ax$ : $f$ est \textbf{linéaire} (elle traduit une situation
        de proportionnalité).
  \item Si $a=0$, $f(x)=b$ : $f$ est \textbf{constante}.
\end{itemize}
\end{definitionbox}

\begin{proprietebox}[Propriété 1 --- Droite représentative]
La courbe d'une fonction affine $f(x)=ax+b$ est une \textbf{droite} :
\begin{itemize}
  \item $b$ est l'\textbf{ordonnée à l'origine} : la droite coupe l'axe des ordonnées
        au point $(0\,;b)$ ;
  \item $a$ est le \textbf{coefficient directeur} : quand $x$ augmente de $1$, $y$
        augmente de $a$ (ou diminue, si $a$ est négatif).
\end{itemize}
La droite d'une fonction linéaire passe par l'origine du repère.
\end{proprietebox}

\begin{proprietebox}[Propriété 2 --- Variations d'une fonction affine]
Soit $f(x)=ax+b$ une fonction affine.
\begin{itemize}
  \item Si $a>0$, $f$ est croissante sur $\R$.
  \item Si $a<0$, $f$ est décroissante sur $\R$.
  \item Si $a=0$, $f$ est constante sur $\R$.
\end{itemize}
\end{proprietebox}

\textbf{Démonstration}\autableau
\begin{correction}
Soit $u$ et $v$ deux nombres tels que $u<v$ ; on compare $f(u)$ et $f(v)$ en
étudiant le signe de leur différence.
\begin{calculs}
f(v)-f(u) &= (av+b)-(au+b) &\qquad& \rem{définition de $f$}\\
f(v)-f(u) &= av+b-au-b && \rem{on supprime les parenthèses}\\
f(v)-f(u) &= a(v-u) && \rem{on factorise par $a$}
\end{calculs}
\explique{$u<v$, donc $v-u>0$ : $f(v)-f(u)$ a le signe de $a$.}{règle des signes d'un
produit}
\explique{Si $a>0$ : $f(v)-f(u)>0$, donc $f(u)<f(v)$ et $f$ est croissante.}{l'ordre
est conservé}
\explique{Si $a<0$ : $f(v)-f(u)<0$, donc $f(u)>f(v)$ et $f$ est décroissante.}{l'ordre
est inversé}
\explique{Si $a=0$ : $f(v)-f(u)=0$, donc $f(u)=f(v)$ et $f$ est constante.}{}
\end{correction}

\begin{methodebox}[Méthode 5 --- Reconnaître une fonction affine et donner ses variations]
\etapeexemples
Pour chaque fonction $f$, $g$, $h$, dire si elle est affine ; si oui, donner $a$, $b$
et ses variations.
\begin{questions}[label=\alph*)]
  \item ${f(x)=5-3x}$
  \item ${g(x)=(x+2)^{2}-x^{2}}$
  \item ${h(x)=x(x+1)}$
\end{questions}
\begin{correction}
\cas{a)}
\begin{calculs}
f(x) &= 5-3x &\qquad& \rem{on recopie l'expression}\\
f(x) &= -3x+5 && \rem{on écrit d'abord le terme en $x$}
\end{calculs}
\explique{$f$ est affine : $a=-3$ et $b=5$.}{forme $ax+b$}
\explique{$a<0$ : $f$ est décroissante sur $\R$.}{propriété 2}
\cas{b)}
\begin{calculs}
g(x) &= (x+2)^{2}-x^{2} &\qquad& \rem{on recopie l'expression}\\
g(x) &= x^{2}+4x+4-x^{2} && \rem{identité remarquable}\\
g(x) &= 4x+4 && \rem{les $x^{2}$ s'annulent}
\end{calculs}
\explique{$g$ est affine : $a=4$ et $b=4$ ; $a>0$, donc $g$ est croissante sur
$\R$.}{il fallait développer pour le voir}
\cas{c)}
\begin{calculs}
h(x) &= x(x+1) &\qquad& \rem{on recopie l'expression}\\
h(x) &= x^{2}+x && \rem{on développe}
\end{calculs}
\explique{$h$ n'est pas affine.}{il reste un terme en $x^{2}$}
\end{correction}

\etape{À vous de jouer}
Même consigne pour les fonctions $p$, $q$ et $r$ :
\[
  p(x)=2(x-3)+x \qquad q(x)=-\frac{x}{2} \qquad r(x)=(x-1)^{2}-(x+1)^{2}
\]
\reponse{$p(x)=3x-6$, croissante ; $q$ linéaire, $a=-\frac{1}{2}$ et $b=0$,
décroissante ; $r(x)=-4x$, linéaire, décroissante.}
\end{methodebox}

\begin{methodebox}[Méthode 6 --- Tracer la droite d'une fonction affine]
\etapeexemples
Dans un même repère, tracer les droites représentant les fonctions
${f(x)=2x-3}$ et ${g(x)=-\frac{1}{2}x+1}$.
\begin{correction}
\cas{Résolution}\par\nopagebreak
\begin{minipage}[c]{0.58\linewidth}\raggedright
\explique{Deux points suffisent pour tracer une droite.}{on calcule deux images}
\explique{$f(0)=-3$ et $f(3)=6-3=3$.}{points $(0\,;-3)$ et $(3\,;3)$}
\explique{$g(0)=1$ et $g(4)=-2+1=-1$.}{$x=4$ est pair : la moitié est entière}
\explique{On place les points, puis on trace chaque droite à la règle.}{contrôle :
$f$ monte de $2$ quand $x$ avance de $1$}
\end{minipage}\hfill
\begin{minipage}[c]{0.38\linewidth}\centering
\begin{tikzpicture}[x=0.45cm,y=0.45cm]
  \repere{-2}{-4}{5}{4}
  \gradx{-1,1,2,3,4} \grady{-3,-2,-1,1,2,3}
  \draw[coulCourbe,line width=1.1pt] (-0.5,-4)--(3.5,4);
  \draw[coulCourbeB,line width=1.1pt] (-2,2)--(5,-1.5);
  \path[coulCourbe] (0,-3) node[croix]{} (3,3) node[croix]{};
  \path[coulCourbeB] (0,1) node[croix]{} (4,-1) node[croix]{};
  \node[coulCourbe,font=\footnotesize,right] at (3.3,3.4) {$\Cf$};
  \node[coulCourbeB,font=\footnotesize] at (-1.5,3) {$\Cg$};
\end{tikzpicture}
\end{minipage}
\end{correction}

\etape{À vous de jouer}
Tracer les droites représentant ${h(x)=-x+2}$ et ${k(x)=\frac{1}{3}x-1}$.
\reponse{$h$ : points $(0\,;2)$ et $(2\,;0)$ ; $k$ : points $(0\,;-1)$ et $(3\,;0)$.}
\end{methodebox}

\afaire{exercices 15 à 19}

\partiecours{Fonction affine par lecture graphique}

\begin{remarquebox}[Lire le coefficient directeur $a$]
\begin{minipage}[c]{0.6\linewidth}\raggedright
On part d'un point de la droite situé sur un nœud du quadrillage et on va, vers la
droite, jusqu'à un autre nœud de la droite :
\[
  a=\frac{\text{déplacement vertical}}{\text{déplacement horizontal}}
\]
Le déplacement vertical est \textbf{positif} si l'on monte, \textbf{négatif} si l'on
descend. Ci-contre : on avance de $2$ et on monte de $3$, donc $a=\dfrac{3}{2}$.
\end{minipage}\hfill
\begin{minipage}[c]{0.36\linewidth}\centering
\begin{tikzpicture}[x=0.5cm,y=0.5cm]
  \quadrillage{0}{0}{5}{5}
  \draw[coulCourbe,line width=1.1pt] (0.333,0)--(3.667,5);
  \draw[-{Stealth[length=2mm]},thick,coulExo] (1,1)--(3,1)
    node[midway,below,font=\footnotesize]{$+2$};
  \draw[-{Stealth[length=2mm]},thick,coulExo] (3,1)--(3,4)
    node[midway,right,font=\footnotesize]{$+3$};
  \path[coulCourbe] (1,1) node[croix]{} (3,4) node[croix]{};
\end{tikzpicture}
\end{minipage}
\end{remarquebox}

\begin{methodebox}[Méthode 7 --- Retrouver l'expression d'une fonction affine sur sa droite]
\etapeexemples
\begin{minipage}[c]{0.58\linewidth}\raggedright
Les droites $(d_1)$, $(d_2)$ et $(d_3)$ représentent trois fonctions affines $f$, $g$
et $h$. Donner l'expression de chacune.
\end{minipage}\hfill
\begin{minipage}[c]{0.38\linewidth}\centering
\begin{tikzpicture}[x=0.42cm,y=0.42cm]
  \repere{-4}{-4}{5}{4}
  \gradx{-3,-2,-1,1,2,3,4} \grady{-3,-2,-1,1,2,3}
  \draw[coulCourbe,line width=1.1pt] (-1.5,-4)--(2.5,4);
  \draw[coulCourbeB,line width=1.1pt] (-3,4)--(5,-1.333);
  \draw[coulCourbeC,line width=1.1pt] (-4,-3)--(5,-3);
  \node[coulCourbe,font=\footnotesize,right] at (2.4,3.6) {$(d_1)$};
  \node[coulCourbeB,font=\footnotesize,above] at (-2.6,4) {$(d_2)$};
  \node[coulCourbeC,font=\footnotesize,above] at (4,-3) {$(d_3)$};
\end{tikzpicture}
\end{minipage}
\begin{correction}
\cas{$(d_1)$}
\explique{$b=-1$.}{la droite coupe l'axe des ordonnées en $(0\,;-1)$}
\explique{$a=\dfrac{2}{1}=2$.}{de $(0\,;-1)$ à $(1\,;1)$ : on avance de $1$ et on
monte de $2$}
\explique{$f(x)=2x-1$.}{}
\cas{$(d_2)$}
\explique{$b=2$.}{elle coupe l'axe des ordonnées en $(0\,;2)$}
\explique{$a=\dfrac{-2}{3}=-\dfrac{2}{3}$.}{de $(0\,;2)$ à $(3\,;0)$ : on avance de
$3$ et on descend de $2$ (déplacement vertical $-2$)}
\explique{$g(x)=-\dfrac{2}{3}x+2$.}{en avançant de $1$, on ne tombe pas sur un
nœud : on avance de $3$}
\cas{$(d_3)$}
\explique{$h(x)=-3$.}{droite horizontale : $a=0$, fonction constante}
\end{correction}

\etape{À vous de jouer}
\begin{minipage}[c]{0.58\linewidth}\raggedright
Les droites $(d_4)$ et $(d_5)$ représentent deux fonctions affines $k$ et $\ell$.
Donner l'expression de chacune.
\reponse{$k(x)=\frac{3}{2}x$ (linéaire) ; $\ell(x)=-2x+2$.}
\end{minipage}\hfill
\begin{minipage}[c]{0.38\linewidth}\centering
\begin{tikzpicture}[x=0.42cm,y=0.42cm]
  \repere{-4}{-4}{4}{4}
  \gradx{-3,-2,-1,1,2,3} \grady{-3,-2,-1,1,2,3}
  \draw[coulCourbe,line width=1.1pt] (-2.667,-4)--(2.667,4);
  \draw[coulCourbeB,line width=1.1pt] (-1,4)--(3,-4);
  \node[coulCourbe,font=\footnotesize,right] at (2.6,3.6) {$(d_4)$};
  \node[coulCourbeB,font=\footnotesize,left] at (-1,3.6) {$(d_5)$};
\end{tikzpicture}
\end{minipage}
\end{methodebox}

\afaire{exercices 20 à 22}

\partiecours{Fonction affine passant par deux points}

\begin{proprietebox}[Propriété 3 --- Coefficient directeur]
Si $A(x_{A}\,;y_{A})$ et $B(x_{B}\,;y_{B})$ sont deux points de la droite d'une fonction
affine $f(x)=ax+b$, avec $x_{A}\neq x_{B}$, alors :
\[
  a=\frac{y_{B}-y_{A}}{x_{B}-x_{A}}
  \qquad\text{(différence des ordonnées sur différence des abscisses)}
\]
\end{proprietebox}

\textbf{Démonstration}\autableau
\begin{correction}
\begin{calculs}
y_{B}-y_{A} &= f(x_{B})-f(x_{A}) &\qquad& \rem{$A$ et $B$ sont sur la droite}\\
y_{B}-y_{A} &= (ax_{B}+b)-(ax_{A}+b) && \rem{définition de $f$}\\
y_{B}-y_{A} &= a(x_{B}-x_{A}) && \rem{les $b$ s'annulent ; on factorise par $a$}
\end{calculs}
\explique{On divise par $x_{B}-x_{A}$, qui n'est pas nul.}{on obtient la formule}
\end{correction}

\begin{methodebox}[Méthode 8 --- Trouver la fonction affine qui passe par deux points]
\etapeexemples
\begin{questions}[label=\alph*)]
  \item Trouver la fonction affine $f$ telle que $f(1)=5$ et $f(3)=-1$.
  \item Une course de taxi de $4$~km coûte $11$~€ ; une course de $10$~km coûte
        $20$~€. Le prix est une fonction affine $p$ de la distance $x$ (en km).
        Trouver $p(x)$.
\end{questions}
\begin{correction}
\cas{a)}
\explique{La droite passe par $A(1\,;5)$ et $B(3\,;-1)$.}{$f(1)=5$ : point
d'abscisse $1$ et d'ordonnée $5$}
\begin{calculs}
a &= \frac{y_{B}-y_{A}}{x_{B}-x_{A}} &\qquad& \rem{propriété 3}\\
a &= \frac{-1-5}{3-1} && \rem{on remplace}\\
a &= -3 && \rem{$-6\div 2$}
\end{calculs}
\explique{Donc $f(x)=-3x+b$ ; on cherche $b$ avec le point $A$.}{$f(1)=5$}
\begin{calculs}
f(1) &= -3\times 1+b &\qquad& \rem{on remplace $x$ par $1$}\\
5 &= -3+b && \rem{car $f(1)=5$}\\
b &= 5+3 && \rem{on ajoute $3$ des deux côtés}\\
b &= 8 &&
\end{calculs}
\explique{$f(x)=-3x+8$.}{contrôle avec $B$ : $f(3)=-9+8=-1$}
\cas{b)}
\explique{La droite passe par $A(4\,;11)$ et $B(10\,;20)$.}{distance en abscisse,
prix en ordonnée}
\begin{calculs}
a &= \frac{20-11}{10-4} &\qquad& \rem{propriété 3}\\
a &= \frac{9}{6} && \\
a &= 1{,}5 && \rem{le prix augmente de $1{,}5$~€ par km}
\end{calculs}
\begin{calculs}
p(4) &= 1{,}5\times 4+b &\qquad& \rem{on remplace $x$ par $4$}\\
11 &= 6+b && \rem{car $p(4)=11$}\\
b &= 5 && \rem{on enlève $6$ des deux côtés}
\end{calculs}
\explique{$p(x)=1{,}5x+5$.}{$5$~€ de prise en charge, plus $1{,}5$~€ par km}
\end{correction}

\etape{À vous de jouer}
\begin{questions}[label=\alph*)]
  \item Trouver la fonction affine $f$ telle que $f(-1)=6$ et $f(3)=-2$.
  \item Trouver la fonction affine $g$ dont la droite passe par $A(-2\,;0)$ et $B(4\,;3)$.
\end{questions}
\reponse{a) $f(x)=-2x+4$ ; b) $g(x)=\frac{1}{2}x+1$.}
\end{methodebox}

\afaire{exercices 23 à 26}

\partiecours{Signe d'une fonction par lecture graphique}

\begin{definitionbox}[Définition 5 --- Signe d'une fonction, tableau de signes]
Étudier le \textbf{signe} d'une fonction $f$, c'est trouver pour quels $x$ :
\begin{itemize}
  \item $f(x)>0$ : la courbe est \textbf{au-dessus} de l'axe des abscisses ;
  \item $f(x)<0$ : la courbe est \textbf{au-dessous} de l'axe des abscisses ;
  \item $f(x)=0$ : la courbe \textbf{coupe} l'axe des abscisses.
\end{itemize}
On le résume dans un \textbf{tableau de signes} : les valeurs de $x$ sur la première
ligne, les signes $+$ ou $-$ et les $0$ sur la seconde.
\end{definitionbox}

\begin{methodebox}[Méthode 9 --- Lire le signe d'une fonction sur sa courbe]
\etapeexemples
\begin{minipage}[c]{0.58\linewidth}\raggedright
La courbe $\Cf$ représente une fonction $f$ définie sur $\intff{-4}{4}$ ; la droite
$\Cg$ représente une fonction affine $g$ définie sur $\R$. Dresser le tableau de
signes de $f$, puis celui de $g$.
\end{minipage}\hfill
\begin{minipage}[c]{0.38\linewidth}\centering
\begin{tikzpicture}[x=0.42cm,y=0.42cm]
  \repere{-5}{-4}{5}{4}
  \gradx{-4,-3,-2,-1,1,2,3,4} \grady{-3,-2,-1,1,2,3}
  \draw[coulCourbe,line width=1.1pt]
    (-4,-2) .. controls (-3.67,-1.33) and (-3.33,-0.56) .. (-3,0)
            .. controls (-2.33,1.11) and (-1.67,3) .. (-1,3)
            .. controls (-0.33,3) and (0.33,0.89) .. (1,0)
            .. controls (1.33,-0.44) and (1.67,-1) .. (2,-1)
            .. controls (2.33,-1) and (2.67,-0.5) .. (3,0)
            .. controls (3.33,0.5) and (3.67,1.33) .. (4,2);
  \path[coulCourbe] (-4,-2) node[croix]{} (4,2) node[croix]{};
  \draw[coulCourbeB,line width=1.1pt] (0,4)--(4,-4);
  \node[coulCourbe,font=\footnotesize] at (-3.2,2.4) {$\Cf$};
  \node[coulCourbeB,font=\footnotesize,right] at (0.6,3.6) {$\Cg$};
\end{tikzpicture}
\end{minipage}
\begin{correction}
\cas{Signe de $f$}
\explique{$\Cf$ coupe l'axe des abscisses en $x=-3$, en $x=1$ et en $x=3$.}{là,
$f(x)=0$}
\explique{De gauche à droite : au-dessous, au-dessus, au-dessous, au-dessus.}{signes
$-$ puis $+$ puis $-$ puis $+$}
\begin{center}
\begin{tikzpicture}
  \tkzTabInit[lgt=2.2,espcl=1.6]{$x$/1,$f(x)$/1}{$-4$,$-3$,$1$,$3$,$4$}
  \tkzTabLine{,-,z,+,z,-,z,+,}
\end{tikzpicture}
\end{center}
\cas{Signe de $g$}
\explique{$\Cg$ coupe l'axe en $x=2$ ; au-dessus avant, au-dessous après.}{la droite
descend}
\begin{center}
\begin{tikzpicture}
  \tkzTabInit[lgt=2.2,espcl=2.2]{$x$/1,$g(x)$/1}{$-\infty$,$2$,$+\infty$}
  \tkzTabLine{,+,z,-,}
\end{tikzpicture}
\end{center}
\end{correction}

\etape{À vous de jouer}
Dresser le tableau de signes des fonctions $k$ et $\ell$ de la méthode 7 (droites
$(d_4)$ et $(d_5)$).
\reponse{$k$ : $-$ avant $0$, $+$ après ; $\ell$ : $+$ avant $1$, $-$ après.}
\end{methodebox}

\afaire{exercices 27 à 29}

\partiecours{Propriétés des inégalités, encadrer une quantité}

\begin{proprietebox}[Propriété 4 --- Opérations sur les inégalités]
\begin{minipage}[c]{0.58\linewidth}\raggedright
Soit $a<b$, et $c$ un nombre.
\begin{itemize}
  \item Ajouter ou soustraire $c$ \textbf{conserve} le sens :
        $a+c<b+c$ et $a-c<b-c$.\\
        \textit{Exemple :} $2<5$, donc $2+3<5+3$, soit $5<8$.
  \item Multiplier ou diviser par $c>0$ \textbf{conserve} le sens :
        $a\times c<b\times c$.\\
        \textit{Exemple :} $2<5$, donc $2\times 4<5\times 4$, soit $8<20$.
  \item Multiplier ou diviser par $c<0$ \textbf{change} le sens :
        $a\times c>b\times c$.\\
        \textit{Exemple :} $2<5$, donc $2\times(-3)>5\times(-3)$, soit $-6>-15$.
\end{itemize}
\end{minipage}\hfill
\begin{minipage}[c]{0.38\linewidth}\centering
\begin{tikzpicture}[x=0.4cm,y=0.4cm]
  \draw[-{Stealth[length=2mm]},thick] (-6.5,0)--(6.5,0);
  \foreach \t in {-5,...,5} \draw (\t,0.15)--(\t,-0.15);
  \node[below,font=\small] at (0,-0.15) {$0$};
  \path[coulCourbe] (2,0) node[croix]{} (5,0) node[croix]{};
  \path[coulCourbeB] (-2,0) node[croix]{} (-5,0) node[croix]{};
  \node[coulCourbe,above,font=\small] at (2,0.2) {$2$};
  \node[coulCourbe,above,font=\small] at (5,0.2) {$5$};
  \node[coulCourbeB,above,font=\small] at (-2,0.2) {$-2$};
  \node[coulCourbeB,above,font=\small] at (-5,0.2) {$-5$};
\end{tikzpicture}

{\small $2<5$, mais $-2>-5$ :\\ multiplier par $-1$ renverse l'ordre}
\end{minipage}
\end{proprietebox}

\begin{methodebox}[Méthode 10 --- Encadrer une quantité]
\etapeexemples
On sait que $2\le x\le 5$. Encadrer chaque expression $A$, $B$.
\begin{questions}[label=\alph*)]
  \item $A=3x-1$
  \item $B=-2x+4$
\end{questions}
\begin{correction}
\cas{a)}
\begin{calculs}
\renewcommand{\arraystretch}{1.5}
\begin{array}{@{}r@{\;}c@{\;}c@{\;}c@{\;}l@{\qquad}l@{}}
2 & \le & x & \le & 5 & \rem{on recopie l'encadrement}\\
6 & \le & 3x & \le & 15 & \rem{on multiplie par $3>0$ : même sens}\\
5 & \le & 3x-1 & \le & 14 & \rem{on enlève $1$}
\end{array}
\end{calculs}
\explique{$5\le A\le 14$.}{conclusion}
\cas{b)}
\begin{calculs}
\renewcommand{\arraystretch}{1.5}
\begin{array}{@{}r@{\;}c@{\;}c@{\;}c@{\;}l@{\qquad}l@{}}
2 & \le & x & \le & 5 & \rem{on recopie l'encadrement}\\
-4 & \ge & -2x & \ge & -10 & \rem{on multiplie par $-2<0$ : le sens change}\\
0 & \ge & -2x+4 & \ge & -6 & \rem{on ajoute $4$}
\end{array}
\end{calculs}
\explique{$-6\le B\le 0$.}{on réécrit l'encadrement dans l'ordre croissant}
\end{correction}

\etape{À vous de jouer}
On sait que $-1\le x\le 3$. Encadrer $C=4x+2$, puis $D=5-x$.
\reponse{$-2\le C\le 14$ ; $2\le D\le 6$.}
\end{methodebox}

\afaire{exercices 30 à 32}

\partiecours{Résoudre une inéquation du premier degré}

Une inéquation du \textbf{premier degré} ne contient que des termes en $x$ et des
nombres. On la résout comme une équation, avec la propriété 4 : les $x$ d'un côté,
les nombres de l'autre, puis on divise par le coefficient de $x$, \textbf{en changeant
le sens s'il est négatif}. On conclut par un intervalle.

\smallskip
\textbf{Changer le sens}, c'est remplacer $<$ par $>$, $>$ par $<$, $\le$ par $\ge$ et
$\ge$ par $\le$. \textit{Exemple :} $-2x<6$ devient $x>-3$.

\begin{methodebox}[Méthode 11 --- Résoudre une inéquation du premier degré]
\etapeexemples
Résoudre chaque inéquation.
\begin{questions}[label=\alph*)]
  \item $3x-7<8$
  \item $5-2x\ge 11$
  \item $4x+1>7x-8$
\end{questions}
\begin{correction}
\cas{a)}
\begin{calculs}
3x-7 &< 8 &\qquad& \rem{on recopie l'inéquation}\\
3x &< 15 && \rem{on ajoute $7$}\\
x &< 5 && \rem{on divise par $3>0$ : même sens}
\end{calculs}
\explique{$S=\intoo{-\infty}{5}$.}{inégalité stricte : $5$ exclu}
\cas{b)}
\begin{calculs}
5-2x &\ge 11 &\qquad& \rem{on recopie l'inéquation}\\
-2x &\ge 6 && \rem{on enlève $5$}\\
x &\le -3 && \rem{on divise par $-2<0$ : le sens change}
\end{calculs}
\explique{$S=\intof{-\infty}{-3}$.}{« $\le$ » : $-3$ compris}
\cas{c)}
\begin{calculs}
4x+1 &> 7x-8 &\qquad& \rem{on recopie l'inéquation}\\
4x-7x &> -8-1 && \rem{les $x$ à gauche, les nombres à droite}\\
-3x &> -9 && \rem{on réduit}\\
x &< 3 && \rem{on divise par $-3<0$ : le sens change}
\end{calculs}
\explique{$S=\intoo{-\infty}{3}$.}{contrôle avec $x=0$ : $1>-8$ est vrai}
\end{correction}

\etape{À vous de jouer}
Résoudre $2x+9\le 1$, puis $6-3x>0$, puis $5x-4\ge 2x+11$.
\reponse{$S=\intof{-\infty}{-4}$ ; $S=\intoo{-\infty}{2}$ ; $S=\intfo{5}{+\infty}$.}
\end{methodebox}

\afaire{exercices 33 à 36}

\besoinplace{9cm}
\partiecours{Signe d'une fonction affine par le calcul}

\begin{proprietebox}[Propriété 5 --- Signe de $ax+b$]
Soit $f(x)=ax+b$ avec $a\neq 0$. On trouve d'abord la valeur qui annule $f(x)$, en
résolvant $ax+b=0$ ; le signe de $a$ donne ensuite l'ordre des signes :

\smallskip
\begin{minipage}{0.48\linewidth}\centering
{\small $a>0$ : droite qui monte}\\[2pt]
\begin{tikzpicture}[scale=0.9]
  \tkzTabInit[lgt=2.2,espcl=2.2]{$x$/1,$ax+b$/1}{$-\infty$,$-\frac{b}{a}$,$+\infty$}
  \tkzTabLine{,-,z,+,}
\end{tikzpicture}

\smallskip
\begin{tikzpicture}[x=0.8cm,y=0.8cm]
  \draw[-{Stealth[length=2mm]},thick] (-0.2,0)--(4.4,0) node[below left=-1pt]{$x$};
  \draw[coulCourbe,line width=1.1pt] (0.6,-1)--(3.4,1);
  \path[coulCourbe] (2,0) node[croix]{};
  \node[below right,font=\small] at (2,0) {$-\frac{b}{a}$};
  \node[coulCourbe,font=\small,anchor=west] at (1,-1.15) {$-$ : au-dessous};
  \node[coulCourbe,font=\small,anchor=east] at (3,1.15) {$+$ : au-dessus};
\end{tikzpicture}
\end{minipage}\hfill
\begin{minipage}{0.48\linewidth}\centering
{\small $a<0$ : droite qui descend}\\[2pt]
\begin{tikzpicture}[scale=0.9]
  \tkzTabInit[lgt=2.2,espcl=2.2]{$x$/1,$ax+b$/1}{$-\infty$,$-\frac{b}{a}$,$+\infty$}
  \tkzTabLine{,+,z,-,}
\end{tikzpicture}

\smallskip
\begin{tikzpicture}[x=0.8cm,y=0.8cm]
  \draw[-{Stealth[length=2mm]},thick] (-0.2,0)--(4.4,0) node[below left=-1pt]{$x$};
  \draw[coulCourbeB,line width=1.1pt] (0.6,1)--(3.4,-1);
  \path[coulCourbeB] (2,0) node[croix]{};
  \node[below left,font=\small] at (2,0) {$-\frac{b}{a}$};
  \node[coulCourbeB,font=\small,anchor=west] at (1,1.15) {$+$ : au-dessus};
  \node[coulCourbeB,font=\small,anchor=east] at (3,-1.15) {$-$ : au-dessous};
\end{tikzpicture}
\end{minipage}

\smallskip
\textit{Pourquoi ?} Pour $a>0$, on résout $ax+b>0$ (méthode 11) :
\begin{calculs}
ax+b &> 0 &\qquad& \rem{$f(x)$ strictement positif}\\
ax &> -b && \rem{on enlève $b$}\\
x &> -\frac{b}{a} && \rem{on divise par $a>0$ : même sens}
\end{calculs}
Donc $+$ après $-\frac{b}{a}$. Pour $a<0$, la division change le sens : $+$ avant
$-\frac{b}{a}$.
\end{proprietebox}

\begin{methodebox}[Méthode 12 --- Dresser le tableau de signes de $ax+b$]
\etapeexemples
Dresser le tableau de signes de chaque expression $A$, $B$.
\begin{questions}[label=\alph*)]
  \item $A(x)=3x-12$
  \item $B(x)=-2x+5$
\end{questions}
\begin{correction}
\cas{a)}
\begin{calculs}
3x-12 &= 0 &\qquad& \rem{valeur qui annule $A(x)$}\\
3x &= 12 && \rem{on ajoute $12$}\\
x &= 4 && \rem{on divise par $3$}
\end{calculs}
\explique{$a=3>0$ : $-$ avant $4$, $+$ après.}{propriété 5 : la droite monte}
\begin{center}
\begin{tikzpicture}
  \tkzTabInit[lgt=2.2,espcl=2.2]{$x$/1,$3x-12$/1}{$-\infty$,$4$,$+\infty$}
  \tkzTabLine{,-,z,+,}
\end{tikzpicture}
\end{center}
\cas{b)}
\begin{calculs}
-2x+5 &= 0 &\qquad& \rem{valeur qui annule $B(x)$}\\
-2x &= -5 && \rem{on enlève $5$}\\
x &= \frac{5}{2} && \rem{on divise par $-2$}
\end{calculs}
\explique{$a=-2<0$ : $+$ avant $\frac{5}{2}$, $-$ après.}{la droite descend}
\begin{center}
\begin{tikzpicture}
  \tkzTabInit[lgt=2.2,espcl=2.2]{$x$/1,$-2x+5$/1}{$-\infty$,$\frac{5}{2}$,$+\infty$}
  \tkzTabLine{,+,z,-,}
\end{tikzpicture}
\end{center}
\end{correction}

\etape{À vous de jouer}
Dresser le tableau de signes de chaque expression $C$, $D$, $E$ :
\[
  C(x)=5x+10 \qquad D(x)=6-4x \qquad E(x)=\frac{1}{2}x-3
\]
\reponse{$C$ : $-$ avant $-2$, $+$ après ; $D$ : $+$ avant $\frac{3}{2}$, $-$ après ;
$E$ : $-$ avant $6$, $+$ après.}
\end{methodebox}

\afaire{exercices 37 à 40}

\partiecours{Inéquations produit}

\begin{proprietebox}[Propriété 6 --- Signe d'un produit]
Le signe d'un produit suit la \textbf{règle des signes} : deux facteurs de même signe
donnent un produit positif, deux facteurs de signes contraires un produit négatif.

Dans un tableau de signes, on met une ligne par facteur, puis une dernière ligne
pour le produit ; un $0$ dans une colonne donne un $0$ au produit.
\end{proprietebox}

\begin{methodebox}[Méthode 13 --- Résoudre une inéquation produit]
\etapeexemples
Résoudre chaque inéquation.
\begin{questions}[label=\alph*)]
  \item $(x-3)(2x+4)<0$
  \item $(5-x)(3x+3)\ge 0$
\end{questions}
\begin{correction}
\cas{a)}
\explique{On cherche les $x$ pour lesquels le produit $(x-3)(2x+4)$ est
\textbf{négatif} : on étudie son signe.}{« $<0$ » veut dire « strictement négatif »}
\explique{$x-3=0$ pour $x=3$ ; $a=1>0$ : $-$ puis $+$.}{signe du premier facteur}
\begin{calculs}
2x+4 &= 0 &\qquad& \rem{second facteur}\\
2x &= -4 &&\\
x &= -2 && \rem{$a=2>0$ : $-$ puis $+$}
\end{calculs}
\explique{Valeurs rangées dans l'ordre : $-2$ puis $3$.}{première ligne du tableau}
\begin{center}
\begin{tikzpicture}
  \tkzTabInit[lgt=3.4,espcl=2]{$x$/1,$x-3$/1,$2x+4$/1,$(x-3)(2x+4)$/1}%
    {$-\infty$,$-2$,$3$,$+\infty$}
  \tkzTabLine{,-,t,-,z,+,}
  \tkzTabLine{,-,z,+,t,+,}
  \tkzTabLine{,+,z,-,z,+,}
\end{tikzpicture}
\end{center}
\explique{$S=\intoo{-2}{3}$.}{on lit les signes $-$ ; inégalité stricte : $-2$ et $3$
exclus}
\cas{b)}
\explique{On cherche les $x$ pour lesquels le produit est \textbf{positif ou
nul}.}{« $\ge 0$ » veut dire « positif ou nul »}
\begin{calculs}
5-x &= 0 &\qquad& \rem{premier facteur}\\
x &= 5 && \rem{$a=-1<0$ : $+$ puis $-$}
\end{calculs}
\begin{calculs}
3x+3 &= 0 &\qquad& \rem{second facteur}\\
3x &= -3 &&\\
x &= -1 && \rem{$a=3>0$ : $-$ puis $+$}
\end{calculs}
\begin{center}
\begin{tikzpicture}
  \tkzTabInit[lgt=3.4,espcl=2]{$x$/1,$5-x$/1,$3x+3$/1,$(5-x)(3x+3)$/1}%
    {$-\infty$,$-1$,$5$,$+\infty$}
  \tkzTabLine{,+,t,+,z,-,}
  \tkzTabLine{,-,z,+,t,+,}
  \tkzTabLine{,-,z,+,z,-,}
\end{tikzpicture}
\end{center}
\explique{$S=\intff{-1}{5}$.}{on lit les signes $+$ ; « $\ge$ » : les $0$ sont
compris}
\end{correction}

\etape{À vous de jouer}
Résoudre $(2x-6)(x+1)\le 0$, puis $x(4-x)>0$.
\reponse{$S=\intff{-1}{3}$ ; $S=\intoo{0}{4}$.}
\end{methodebox}

\afaire{exercices 41 à 43}

\partiecours{Inéquations quotient}

\begin{proprietebox}[Propriété 7 --- Signe d'un quotient]
Un quotient suit la même règle des signes qu'un produit. Le dénominateur ne peut pas
être nul : la valeur qui l'annule est une \textbf{valeur interdite}. Dans le tableau,
on la marque d'une \textbf{double barre}, et elle n'est jamais solution.
\end{proprietebox}

\begin{methodebox}[Méthode 14 --- Résoudre une inéquation quotient]
\etapeexemples
Résoudre chaque inéquation.
\begin{questions}[label=\alph*)]
  \item $\dfrac{3x-6}{x+1}\ge 0$
  \item $\dfrac{4-2x}{x+3}>0$
\end{questions}
\begin{correction}
\cas{a)}
\explique{Valeur interdite : $x+1=0$ pour $x=-1$.}{dénominateur nul}
\explique{$3x-6=0$ pour $x=2$ ; $a=3>0$ : $-$ puis $+$.}{$3x=6$}
\explique{$x+1$ : $a=1>0$ : $-$ puis $+$.}{signe du dénominateur}
\begin{center}
\begin{tikzpicture}
  \tkzTabInit[lgt=3.4,espcl=2]{$x$/1,$3x-6$/1,$x+1$/1,$\dfrac{3x-6}{x+1}$/1.4}%
    {$-\infty$,$-1$,$2$,$+\infty$}
  \tkzTabLine{,-,t,-,z,+,}
  \tkzTabLine{,-,z,+,t,+,}
  \tkzTabLine{,+,d,-,z,+,}
\end{tikzpicture}
\end{center}
\explique{$S=\intoo{-\infty}{-1}\cup\intfo{2}{+\infty}$.}{signes $+$ ; $2$ compris
($\ge$) ; $-1$ exclu (interdite)}
\cas{b)}
\explique{Valeur interdite : $x+3=0$ pour $x=-3$.}{dénominateur nul}
\explique{$4-2x=0$ pour $x=2$ ; $a=-2<0$ : $+$ puis $-$.}{$-2x=-4$}
\begin{center}
\begin{tikzpicture}
  \tkzTabInit[lgt=3.4,espcl=2]{$x$/1,$4-2x$/1,$x+3$/1,$\dfrac{4-2x}{x+3}$/1.4}%
    {$-\infty$,$-3$,$2$,$+\infty$}
  \tkzTabLine{,+,t,+,z,-,}
  \tkzTabLine{,-,z,+,t,+,}
  \tkzTabLine{,-,d,+,z,-,}
\end{tikzpicture}
\end{center}
\explique{$S=\intoo{-3}{2}$.}{signes $+$ ; inégalité stricte}
\end{correction}

\etape{À vous de jouer}
Résoudre $\dfrac{x-5}{x+2}\le 0$, puis $\dfrac{2x+1}{3-x}\ge 0$.
\reponse{$S=\intof{-2}{5}$ ; $S=\intfo{-\frac{1}{2}}{3}$.}
\end{methodebox}

\afaire{exercices 44 à 46}

\partiecours{Inéquations avancées}

Le tableau de signes ne sert qu'à comparer un produit ou un quotient \textbf{à $0$}.
Sinon, on s'y ramène : on passe tout dans le membre de gauche, puis on factorise ou on
met au même dénominateur.

\begin{methodebox}[Méthode 15 --- Se ramener à une comparaison avec $0$]
\etapeexemples
Résoudre chaque inéquation.
\begin{questions}[label=\alph*)]
  \item $x^{2}\ge 5x$
  \item $\dfrac{x+4}{x-2}\le 3$
\end{questions}
\begin{correction}
\cas{a)}
\begin{calculs}
x^{2} &\ge 5x &\qquad& \rem{on recopie l'inéquation}\\
x^{2}-5x &\ge 0 && \rem{on passe $5x$ à gauche}\\
x(x-5) &\ge 0 && \rem{facteur commun $x$}
\end{calculs}
\explique{$x$ s'annule en $0$, $x-5$ en $5$ ; tous deux : $-$ puis $+$.}{coefficients
de $x$ positifs}
\begin{center}
\begin{tikzpicture}
  \tkzTabInit[lgt=3.4,espcl=2]{$x$/1,$x$/1,$x-5$/1,$x(x-5)$/1}%
    {$-\infty$,$0$,$5$,$+\infty$}
  \tkzTabLine{,-,z,+,t,+,}
  \tkzTabLine{,-,t,-,z,+,}
  \tkzTabLine{,+,z,-,z,+,}
\end{tikzpicture}
\end{center}
\explique{$S=\intof{-\infty}{0}\cup\intfo{5}{+\infty}$.}{signes $+$ et les $0$ ;
diviser par $x$ aurait perdu des solutions}
\cas{b)}
\explique{Valeur interdite : $x=2$.}{$x-2=0$}
\begin{calculs}
\frac{x+4}{x-2}-3 &\le 0 &\qquad& \rem{on passe $3$ à gauche}\\
\frac{x+4}{x-2}-\frac{3\fois{(x-2)}}{1\fois{(x-2)}} &\le 0 && \rem{même dénominateur}\\
\frac{x+4-3x+6}{x-2} &\le 0 && \rem{attention au signe $-$ devant la fraction}\\
\frac{-2x+10}{x-2} &\le 0 && \rem{on réduit}
\end{calculs}
\explique{$-2x+10=0$ pour $x=5$ ; $a=-2<0$ : $+$ puis $-$.}{$x-2$ : $-$ puis $+$}
\begin{center}
\begin{tikzpicture}
  \tkzTabInit[lgt=3.4,espcl=2]{$x$/1,$-2x+10$/1,$x-2$/1,$\dfrac{-2x+10}{x-2}$/1.4}%
    {$-\infty$,$2$,$5$,$+\infty$}
  \tkzTabLine{,+,t,+,z,-,}
  \tkzTabLine{,-,z,+,t,+,}
  \tkzTabLine{,-,d,+,z,-,}
\end{tikzpicture}
\end{center}
\explique{$S=\intoo{-\infty}{2}\cup\intfo{5}{+\infty}$.}{signes $-$ ; $5$ compris ;
$2$ interdit}
\end{correction}

\etape{À vous de jouer}
Résoudre $(x+2)^{2}<9$, puis $\dfrac{6}{x+1}\ge 2$.
\reponse{$(x-1)(x+5)<0$, $S=\intoo{-5}{1}$ ; $\dfrac{4-2x}{x+1}\ge 0$,
$S=\intof{-1}{2}$.}
\end{methodebox}

\afaire{exercices 47 à 50}

\partiecours{Problèmes d'inéquations}

\begin{methodebox}[Méthode 16 --- Résoudre un problème avec une inéquation]
\etapeexemples
Pour une fête, la salle A coûte $200$~€ plus $15$~€ par invité ; la salle B coûte
$50$~€ plus $20$~€ par invité. On note $x$ le nombre d'invités.
\begin{questions}[label=\alph*)]
  \item Exprimer les prix $A(x)$ et $B(x)$ des deux salles.
  \item Traduire « la salle A est moins chère que la salle B » par une inéquation.
  \item Résoudre cette inéquation.
  \item Conclure par une phrase.
\end{questions}
\begin{correction}
\cas{a)}
\explique{$A(x)=15x+200$ et $B(x)=20x+50$.}{prix par invité fois $x$, plus le prix
fixe}
\cas{b)}
\explique{$A(x)<B(x)$, soit $15x+200<20x+50$.}{« moins cher » : strictement inférieur}
\cas{c)}
\begin{calculs}
15x+200 &< 20x+50 &\qquad& \rem{on recopie l'inéquation}\\
15x-20x &< 50-200 && \rem{les $x$ à gauche, les nombres à droite}\\
-5x &< -150 && \rem{on réduit}\\
x &> 30 && \rem{on divise par $-5<0$ : le sens change}
\end{calculs}
\cas{d)}
\explique{La salle A est moins chère à partir de $31$ invités.}{$x$ est un nombre
entier ; pour $30$ invités, les prix sont égaux}
\end{correction}

\etape{À vous de jouer}
Le taxi A facture $4$~€ de prise en charge, puis $1{,}50$~€ par kilomètre ; le taxi B
facture $10$~€ de prise en charge, puis $1$~€ par kilomètre. On note $x$ la distance
parcourue, en km.
\begin{questions}
  \item Exprimer les prix $A(x)$ et $B(x)$ des deux taxis.
  \item Traduire « le taxi A est moins cher que le taxi B » par une inéquation.
  \item Résoudre cette inéquation.
  \item Conclure par une phrase.
\end{questions}
\reponse{1\textdegree) $A(x)=1{,}5x+4$ et $B(x)=x+10$ ;\; 2\textdegree)
$1{,}5x+4<x+10$ ;\; 3\textdegree) $x<12$ ;\; 4\textdegree) le taxi A est moins cher pour
moins de $12$~km.}
\end{methodebox}

\afaire{exercices 51 à 53}

\end{document}
